% !TEX root = ../algebra_02.tex

\section{Изменение матрицы линейного отображения при замене базисов}

\begin{Definition}{Матрица перехода}{}
	Пусть $E,E'$ --- два базиса одного пространства $V$. Матрицей перехода от $E'$ к $E$ называется матрица тождественного оператора:
	\[
	C_{E'\to E}:=[\operatorname{id}_V]_{E',E}
	\].
	Она переводит координаты вектора в базисе $E'$ в координаты в базисе $E$, то есть $[v]_E=C_{E'\to E}[v]_{E'}$.
\end{Definition}

\begin{Theorem}{Формула замены базисов}{}
	Пусть $\alpha\colon V\to W$ --- линейное отображение, $E,E'$ --- базисы $V$, а $F,F'$ --- базисы $W$. Тогда матрица отображения в новых базисах выражается формулой:
	\[
	[\alpha]_{E',F'} = C_{F\to F'}[\alpha]_{E,F}C_{E'\to E}
	\].
\end{Theorem}

\begin{proof}
	Представим отображение $\alpha$ в виде композиции с тождественными отображениями: $\alpha = \operatorname{id}_W\circ \alpha\circ \operatorname{id}_V$.
	Используя теорему о матрице композиции, получаем:
	\[
	[\alpha]_{E',F'} = [\operatorname{id}_W]_{F,F'} [\alpha]_{E,F} [\operatorname{id}_V]_{E',E} = C_{F\to F'}[\alpha]_{E,F}C_{E'\to E}
	\].
\end{proof}

В случае линейного оператора (когда $V=W$ и базисы совпадают) формула упрощается.

\begin{Corollary}{Замена базиса для оператора}{}
	Если $\alpha\in \End V$, $E,E'$ --- базисы $V$, $[\alpha]_E=A$ и $C=C_{E'\to E}$. Тогда
	\[
	[\alpha]_{E'}=C^{-1}AC
	\].
\end{Corollary}

Матрицы $A,A'\in M_n(K)$ называются подобными, если существует обратимая матрица $C\in GL_n(K)$ такая, что $A'=C^{-1}AC$. Подобные матрицы имеют геометрический смысл — это матрицы одного и того же линейного оператора в разных базисах.
